% Points-mode starter template for ximeraExam 1.7.1.
% Keep ximeraExam.sty beside this file and compile with LuaLaTeX twice.
\DocumentMetadata{lang=en-US,pdfstandard=ua-2,tagging=on}
\documentclass{ximera}
\usepackage{unicode-math}
\usepackage[points]{ximeraExam}

% Edit these details for your course and exam.
\examtitle{Exam Title}
\examcourse{Course Name}
\examdate{Exam Date}
\examversion{Form A}

% Keep this for student copies; switch to \printanswers for the answer key.
\noprintanswers

\begin{document}
\begin{examcover}
  \examcoverheading{\examtitletext}
  \noindent\examcoursetext\quad\examversiontext\quad\examdatetext\par
  \noindent\examnameprompt\ \examnamefield\par\bigskip
  \noindent\textbf{Instructions:} Show your work and justify your answers.
  \par\bigskip
  \gradetable[h][][4]
\end{examcover}

\begin{questions}

% Fixed response space is visible only in the student copy.
\question[5] Solve \(2x+3=11\). Show your work.
\begin{solution}[1.5in]
  Subtract \(3\), then divide by \(2\): \(x=4\).
\end{solution}

% Stretchable response space expands within the available page space.
\question[5] Find the derivative of \(f(x)=x^3-2x\).
\begin{solution}[\stretch{1}]
  \(f'(x)=3x^2-2\).
\end{solution}

% For nested items, the question inherits points from its parts.
\question Consider the function \(g(x)=x^2\).
\begin{parts}
  \part[2] Compute \(g(3)\).
  \begin{solution}[0.75in]
    \(g(3)=9\).
  \end{solution}

  \part[3] Solve \(g(x)=9\).
  \begin{solution}[1in]
    \(x=-3\) or \(x=3\).
  \end{solution}
\end{parts}

\end{questions}
\end{document}
